\documentclass[aspectratio=169, 12pt]{beamer}

%% ====== PACKAGES ======
\usepackage[utf8]{inputenc}
\usepackage[T1]{fontenc}
\usepackage[french]{babel}
\usepackage{helvet}  % Police Arial/Helvetica
\renewcommand{\familydefault}{\sfdefault}  % Sans-serif par défaut
\usepackage{tikz}
\usetikzlibrary{arrows.meta, positioning, calc, shapes.geometric, fit, backgrounds}
\usepackage{listings}
\usepackage{booktabs}
\usepackage{graphicx}
\usepackage{fontawesome5}
\usepackage{hyperref}
\usepackage{appendixnumberbeamer}
\usepackage{pgfpages}

%% ====== MODE NOTES PRÉSENTATEUR ======
%% Les notes s'affichent sur ton écran, pas sur le projecteur
\setbeameroption{show notes on second screen=right}

%% ====== THEME ======
\usetheme{Madrid}
\usecolortheme{default}

% Couleurs UCA (Université Côte d'Azur)
\definecolor{ucablue}{RGB}{0, 82, 155}
\definecolor{ucadark}{RGB}{0, 50, 100}
\definecolor{ucalight}{RGB}{200, 220, 240}
\definecolor{vertok}{RGB}{40, 160, 60}
\definecolor{rougeko}{RGB}{200, 40, 40}
\definecolor{orangewarn}{RGB}{230, 130, 20}
\definecolor{griscode}{RGB}{245, 245, 245}

\setbeamercolor{palette primary}{bg=ucablue, fg=white}
\setbeamercolor{palette secondary}{bg=ucadark, fg=white}
\setbeamercolor{palette tertiary}{bg=ucadark, fg=white}
\setbeamercolor{structure}{fg=ucablue}
\setbeamercolor{title}{fg=white, bg=ucablue}
\setbeamercolor{frametitle}{fg=white, bg=ucablue}
\setbeamercolor{block title}{bg=ucablue, fg=white}
\setbeamercolor{block body}{bg=ucalight, fg=black}
\setbeamercolor{block title alerted}{bg=rougeko, fg=white}
\setbeamercolor{block body alerted}{bg=red!5, fg=black}
\setbeamercolor{block title example}{bg=vertok, fg=white}
\setbeamercolor{block body example}{bg=green!5, fg=black}

%% ====== LISTINGS ======
\lstset{
    basicstyle=\ttfamily\scriptsize,
    backgroundcolor=\color{griscode},
    frame=single,
    rulecolor=\color{gray!40},
    breaklines=true,
    showstringspaces=false,
    keywordstyle=\color{ucablue}\bfseries,
    commentstyle=\color{green!50!black},
    stringstyle=\color{red!60!black},
    tabsize=4
}

%% ====== FOOTLINE avec numéros ======
\setbeamertemplate{navigation symbols}{}
\setbeamertemplate{footline}{
    \leavevmode\hbox{%
        \begin{beamercolorbox}[wd=.33\paperwidth,ht=2.5ex,dp=1ex,center]{palette primary}
            \usebeamerfont{author in head/foot}\insertshortauthor
        \end{beamercolorbox}%
        \begin{beamercolorbox}[wd=.34\paperwidth,ht=2.5ex,dp=1ex,center]{palette secondary}
            \usebeamerfont{title in head/foot}\insertshorttitle
        \end{beamercolorbox}%
        \begin{beamercolorbox}[wd=.33\paperwidth,ht=2.5ex,dp=1ex,right]{palette tertiary}
            \usebeamerfont{date in head/foot}\insertframenumber{} / \inserttotalframenumber\hspace*{2ex}
        \end{beamercolorbox}%
    }
}

%% ====== INFOS ======
\title[TER S2 -- Groupe D]{Génération de Labyrinthes Pac-Man}
\subtitle{UE Travaux d'Études et de Recherche -- Semestre 2}
\author[Groupe D]{Ikram Benchalal \and Nada Zina \and Aya Haddoun}
\institute[UCA]{Université Côte d'Azur\\Master 1 Informatique -- Parcours IA\\[0.3cm]Superviseur : Gilles Menez}
\date{13 mars 2026}

\begin{document}

%% ====================================================================
%% TITRE
%% ====================================================================
\begin{frame}
\titlepage
\end{frame}

\note{
\textbf{SLIDE 1 -- TITRE (30-40 sec)}\\[0.3cm]
Bonjour à tous,\\[0.2cm]
Je suis Ikram, avec mes collègues Nada et Aya.\\[0.2cm]
On vous présente aujourd'hui l'avancement de notre TER : la génération automatique de labyrinthes pour le jeu Pac-Man.\\[0.2cm]
L'objectif du projet c'est de créer un système complet : des algorithmes qui génèrent les labyrinthes, une API déployée dans le cloud pour y accéder, et une interface graphique.\\[0.2cm]
On va vous montrer où on en est, ce qui fonctionne déjà et ce qu'il reste à faire.\\[0.3cm]
\textit{→ Sommaire}
}

%% ====================================================================
%% SOMMAIRE
%% ====================================================================
\begin{frame}{Sommaire}
\tableofcontents
\end{frame}

\note{
\textbf{SLIDE 2 -- SOMMAIRE (20-30 sec)}\\[0.3cm]
Voici le plan de la présentation.\\[0.2cm]
On va commencer par les \textbf{algorithmes de génération}, c'est ma partie.\\[0.2cm]
Ensuite Nada présentera l'\textbf{API} et le \textbf{déploiement}.\\[0.2cm]
Puis on verra les \textbf{tests} et le \textbf{CI/CD} que j'ai mis en place.\\[0.2cm]
Aya parlera de l'\textbf{interface graphique}.\\[0.2cm]
Et on terminera par la \textbf{répartition du travail} dans l'équipe.\\[0.3cm]
\textit{→ Algos}
}

%% ####################################################################
%% PARTIE IKRAM
%% ####################################################################

\section{Algorithmes de génération}

%% --------------------------------------------------------------------
\begin{frame}{Vue d'ensemble : DFS vs Prim}

\begin{columns}[T]
\begin{column}{0.48\textwidth}
\begin{block}{\centering DFS (Ikram)}
\centering
\begin{tikzpicture}[scale=0.32, every node/.style={minimum size=0.32cm}]
    % Grille 7x7 - DFS : couloirs longs
    \fill[black!80] (0,0) rectangle (7,7);
    % Couloirs DFS (longs et sinueux)
    \foreach \x/\y in {1/1, 1/2, 1/3, 1/4, 1/5,
                        2/5, 3/5, 3/4, 3/3, 3/2, 3/1,
                        4/1, 5/1, 5/2, 5/3, 5/4, 5/5,
                        4/3}{
        \fill[white] (\x, \y) rectangle ++(1,1);
    }
\end{tikzpicture}

\small Creuse en profondeur\\
\textbf{Longs couloirs sinueux}
\end{block}
\end{column}

\begin{column}{0.48\textwidth}
\begin{block}{\centering Prim (Nada + Ikram)}
\centering
\begin{tikzpicture}[scale=0.32, every node/.style={minimum size=0.32cm}]
    % Grille 7x7 - Prim : structure étalée
    \fill[black!80] (0,0) rectangle (7,7);
    % Couloirs Prim (plus étalés, plus d'intersections)
    \foreach \x/\y in {1/1, 2/1, 3/1, 3/2, 3/3,
                        1/3, 2/3, 4/3, 5/3,
                        1/4, 1/5, 2/5, 3/5,
                        5/1, 5/2, 5/4, 5/5, 4/5}{
        \fill[white] (\x, \y) rectangle ++(1,1);
    }
\end{tikzpicture}

\small Choisit un mur au hasard\\
\textbf{Structure étalée, ouverte}
\end{block}
\end{column}
\end{columns}

\vspace{0.3cm}
\centering
\small \textcolor{ucablue}{\faArrowRight} ~Les deux génèrent des labyrinthes \textbf{parfaits} (un seul chemin entre chaque paire de cases)

\end{frame}

\note{
\textbf{SLIDE 3 -- DFS vs PRIM (30-40 sec)}\\[0.3cm]
Pour générer les labyrinthes, on a testé deux algorithmes.\\[0.2cm]
À gauche, le \textbf{DFS} : il creuse toujours tout droit jusqu'à être bloqué, puis il revient en arrière. Résultat : des \textbf{longs couloirs sinueux}.\\[0.2cm]
À droite, \textbf{Prim} : il choisit un mur au hasard à chaque étape. Résultat : une structure plus \textbf{étalée} avec plus d'intersections.\\[0.2cm]
Les deux génèrent des labyrinthes \textbf{parfaits}, c'est-à-dire qu'il existe un seul chemin entre deux cases.\\[0.3cm]
\textit{→ On va voir en détail comment fonctionne DFS}
}

%% --------------------------------------------------------------------
\begin{frame}{DFS pas à pas -- Creuse le plus loin possible}

\centering
\begin{tikzpicture}[scale=0.32]
    % Grille 7x7 (indices 0-6, cases impaires = couloirs potentiels : 1, 3, 5)
    
    % === ÉTAPE 1 : Départ ===
    \node[font=\scriptsize\bfseries, ucablue] at (3.5, 8.2) {1. Départ};
    \foreach \x in {0,...,6} \foreach \y in {0,...,6} \fill[black!80] (\x, \y) rectangle ++(1,1);
    \fill[yellow!50] (1,5) rectangle ++(1,1);
    \node[font=\tiny\bfseries] at (1.5,5.5) {S};
    % Voisins possibles marqués
    \draw[thick, orange, dashed] (1.5,5.5) -- (1.5,3.5);
    \draw[thick, orange, dashed] (1.5,5.5) -- (3.5,5.5);
    \node[font=\tiny, orange] at (1.5,3.2) {?};
    \node[font=\tiny, orange] at (3.8,5.5) {?};
    
    % === ÉTAPE 2 : Va au Sud (creuse 2 cases) ===
    \node[font=\scriptsize\bfseries, ucablue] at (11.5, 8.2) {2. Sud $\times$2};
    \foreach \x in {0,...,6} \foreach \y in {0,...,6} \fill[black!80] (8+\x, \y) rectangle ++(1,1);
    \fill[white] (9,5) rectangle ++(1,1);
    \fill[white] (9,4) rectangle ++(1,1);
    \fill[orange!50] (9,3) rectangle ++(1,1);
    \draw[-{Stealth}, ultra thick, red!70!black] (9.5,5.5) -- (9.5,3.5);
    \node[font=\tiny] at (9.5,5.5) {S};
    
    % === ÉTAPE 3 : Continue Sud vers (1,1) ===
    \node[font=\scriptsize\bfseries, ucablue] at (19.5, 8.2) {3. Sud $\times$2};
    \foreach \x in {0,...,6} \foreach \y in {0,...,6} \fill[black!80] (16+\x, \y) rectangle ++(1,1);
    \fill[white] (17,5) rectangle ++(1,1);
    \fill[white] (17,4) rectangle ++(1,1);
    \fill[white] (17,3) rectangle ++(1,1);
    \fill[white] (17,2) rectangle ++(1,1);
    \fill[orange!50] (17,1) rectangle ++(1,1);
    \draw[-{Stealth}, ultra thick, red!70!black] (17.5,3.5) -- (17.5,1.5);
    \node[font=\tiny] at (17.5,5.5) {S};
    % Montrer qu'on peut encore aller à l'Est!
    \draw[thick, orange, dashed] (17.5,1.5) -- (19.5,1.5);
    \node[font=\tiny, orange] at (19.8,1.5) {?};
    
    % === ÉTAPE 4 : Va à l'Est vers (3,1) ===
    \node[font=\scriptsize\bfseries, ucablue] at (3.5, -0.8) {4. Est $\times$2};
    \foreach \x in {0,...,6} \foreach \y in {-8,...,-2} \fill[black!80] (\x, \y) rectangle ++(1,1);
    \fill[white] (1,-3) rectangle ++(1,1);
    \fill[white] (1,-4) rectangle ++(1,1);
    \fill[white] (1,-5) rectangle ++(1,1);
    \fill[white] (1,-6) rectangle ++(1,1);
    \fill[white] (1,-7) rectangle ++(1,1);
    \fill[white] (2,-7) rectangle ++(1,1);
    \fill[orange!50] (3,-7) rectangle ++(1,1);
    \node[font=\tiny] at (1.5,-2.5) {S};
    \draw[-{Stealth}, ultra thick, red!70!black] (1.5,-6.5) -- (3.5,-6.5);
    % Peut aller Est ou Nord!
    \draw[thick, orange, dashed] (3.5,-6.5) -- (5.5,-6.5);
    \draw[thick, orange, dashed] (3.5,-6.5) -- (3.5,-4.5);
    \node[font=\tiny, orange] at (5.8,-6.5) {?};
    \node[font=\tiny, orange] at (3.5,-4.2) {?};
    
    % === ÉTAPE 5 : Continue Est vers (5,1) puis Nord possible ===
    \node[font=\scriptsize\bfseries, ucablue] at (11.5, -0.8) {5. Est $\times$2};
    \foreach \x in {0,...,6} \foreach \y in {-8,...,-2} \fill[black!80] (8+\x, \y) rectangle ++(1,1);
    \fill[white] (9,-3) rectangle ++(1,1);
    \fill[white] (9,-4) rectangle ++(1,1);
    \fill[white] (9,-5) rectangle ++(1,1);
    \fill[white] (9,-6) rectangle ++(1,1);
    \fill[white] (9,-7) rectangle ++(1,1);
    \fill[white] (10,-7) rectangle ++(1,1);
    \fill[white] (11,-7) rectangle ++(1,1);
    \fill[white] (12,-7) rectangle ++(1,1);
    \fill[orange!50] (13,-7) rectangle ++(1,1);
    \node[font=\tiny] at (9.5,-2.5) {S};
    \draw[-{Stealth}, ultra thick, red!70!black] (11.5,-6.5) -- (13.5,-6.5);
    % Peut encore aller au Nord!
    \draw[thick, orange, dashed] (13.5,-6.5) -- (13.5,-4.5);
    \node[font=\tiny, orange] at (13.5,-4.2) {?};
    \node[font=\tiny, rougeko] at (14.7,-6.5) {bord};
    
    % === ÉTAPE 6 : Monte au Nord vers (5,3) ===
    \node[font=\scriptsize\bfseries, ucablue] at (19.5, -0.8) {6. Nord $\times$2};
    \foreach \x in {0,...,6} \foreach \y in {-8,...,-2} \fill[black!80] (16+\x, \y) rectangle ++(1,1);
    % Couloir en L puis remonte
    \fill[white] (17,-3) rectangle ++(1,1);
    \fill[white] (17,-4) rectangle ++(1,1);
    \fill[white] (17,-5) rectangle ++(1,1);
    \fill[white] (17,-6) rectangle ++(1,1);
    \fill[white] (17,-7) rectangle ++(1,1);
    \fill[white] (18,-7) rectangle ++(1,1);
    \fill[white] (19,-7) rectangle ++(1,1);
    \fill[white] (20,-7) rectangle ++(1,1);
    \fill[white] (21,-7) rectangle ++(1,1);
    \fill[white] (21,-6) rectangle ++(1,1);
    \fill[orange!50] (21,-5) rectangle ++(1,1);
    \node[font=\tiny] at (17.5,-2.5) {S};
    \draw[-{Stealth}, ultra thick, red!70!black] (21.5,-6.5) -- (21.5,-4.5);
    % Peut encore aller Nord ou Ouest!
    \draw[thick, orange, dashed] (21.5,-4.5) -- (21.5,-2.5);
    \draw[thick, orange, dashed] (21.5,-4.5) -- (19.5,-4.5);
    \node[font=\tiny, orange] at (21.5,-2.2) {?};
    \node[font=\tiny, orange] at (19.2,-4.5) {?};
    
    % Légende à droite
    \node[right, font=\scriptsize] at (23, -2.5) {\colorbox{yellow!50}{~S~} départ};
    \node[right, font=\scriptsize] at (23, -3.5) {\colorbox{orange!50}{~~} position actuelle};
    \node[right, font=\scriptsize] at (23, -4.5) {\colorbox{red!30}{~~} bloqué};
    \node[right, font=\scriptsize] at (23, -5.5) {\colorbox{white}{~~} couloir creusé};
    \node[right, font=\scriptsize] at (23, -6.5) {\colorbox{black!80}{\textcolor{white}{~~}} mur};
\end{tikzpicture}

\vspace{0.2cm}
\small \textcolor{ucablue}{\faArrowRight} \textbf{Règle :} tant qu'on peut avancer $\rightarrow$ on avance ~\textbar~ \textbf{Sinon} $\rightarrow$ \textcolor{rougeko}{backtrack}

\end{frame}

\note{
\textbf{SLIDE 4 -- DFS PAS À PAS (40-50 sec)}\\[0.3cm]
Je vous montre le principe de DFS sur une grille.\\[0.2cm]
\textbf{Étape 1} : On démarre en S. On regarde les voisins à 2 cases de distance.\\[0.2cm]
\textbf{Étapes 2-3} : On choisit Sud et on \textbf{creuse de 2 cases}. On arrive en (1,1). Sud = bord, mais Est = libre!\\[0.2cm]
\textbf{Étapes 4-5} : On va à l'Est vers (3,1) puis (5,1). Est = bord maintenant, mais Nord = libre!\\[0.2cm]
\textbf{Étape 6} : On monte au Nord vers (5,3). On peut encore continuer Nord ou Ouest...\\[0.2cm]
Le DFS continue \textbf{tant qu'il y a un voisin non visité}. Il ne fait backtrack que quand il est \textbf{coincé de TOUS les côtés}.\\[0.2cm]
C'est pour ça que DFS crée des \textbf{longs couloirs sinueux} : il va toujours le plus loin possible!\\[0.3cm]
\textit{→ Prim fonctionne différemment}
}

%% --------------------------------------------------------------------
\begin{frame}{Prim pas à pas -- Évolution étape par étape}

\centering
\begin{tikzpicture}[scale=0.38]
    % === ÉTAPE 1 : Départ ===
    \node[font=\scriptsize\bfseries, ucablue] at (2.5, 6) {1. Départ};
    \foreach \x in {0,...,4} \foreach \y in {0,...,4} \fill[black!80] (\x, \y) rectangle ++(1,1);
    \fill[yellow!50] (1,3) rectangle ++(1,1);
    \fill[orange!50] (1,1) rectangle ++(1,1);
    \fill[orange!50] (3,3) rectangle ++(1,1);
    \node[font=\tiny\bfseries] at (1.5,3.5) {S};
    \node[font=\tiny, white] at (1.5,1.5) {?};
    \node[font=\tiny, white] at (3.5,3.5) {?};
    
    % === ÉTAPE 2 : Pioche (3,3) ===
    \node[font=\scriptsize\bfseries, ucablue] at (8.5, 6) {2. Pioche};
    \foreach \x in {0,...,4} \foreach \y in {0,...,4} \fill[black!80] (6+\x, \y) rectangle ++(1,1);
    \fill[white] (7,3) rectangle ++(1,1);
    \fill[white] (8,3) rectangle ++(1,1);
    \fill[white] (9,3) rectangle ++(1,1);
    \fill[orange!50] (7,1) rectangle ++(1,1);
    \fill[orange!50] (9,1) rectangle ++(1,1);
    \node[font=\tiny, white] at (7.5,1.5) {?};
    \node[font=\tiny, white] at (9.5,1.5) {?};
    \draw[-{Stealth}, thick, vertok] (7.5,3.5) -- (9.5,3.5);
    
    % === ÉTAPE 3 : Pioche (1,1) ===
    \node[font=\scriptsize\bfseries, ucablue] at (14.5, 6) {3. Pioche};
    \foreach \x in {0,...,4} \foreach \y in {0,...,4} \fill[black!80] (12+\x, \y) rectangle ++(1,1);
    \fill[white] (13,3) rectangle ++(1,1);
    \fill[white] (14,3) rectangle ++(1,1);
    \fill[white] (15,3) rectangle ++(1,1);
    \fill[white] (13,2) rectangle ++(1,1);
    \fill[white] (13,1) rectangle ++(1,1);
    \fill[orange!50] (15,1) rectangle ++(1,1);
    \fill[orange!50] (13,-1) rectangle ++(1,1);
    \node[font=\tiny, white] at (15.5,1.5) {?};
    \draw[-{Stealth}, thick, vertok] (13.5,3.5) -- (13.5,1.5);
    
    % === ÉTAPE 4 : Continue ===
    \node[font=\scriptsize\bfseries, ucablue] at (2.5, -1.5) {4. Pioche};
    \foreach \x in {0,...,4} \foreach \y in {-6,...,-2} \fill[black!80] (\x, \y) rectangle ++(1,1);
    \fill[white] (1,-3) rectangle ++(1,1);
    \fill[white] (2,-3) rectangle ++(1,1);
    \fill[white] (3,-3) rectangle ++(1,1);
    \fill[white] (1,-4) rectangle ++(1,1);
    \fill[white] (1,-5) rectangle ++(1,1);
    \fill[white] (3,-4) rectangle ++(1,1);
    \fill[white] (3,-5) rectangle ++(1,1);
    \fill[orange!50] (1,-7) rectangle ++(1,1);
    \fill[orange!50] (3,-7) rectangle ++(1,1);
    
    % === ÉTAPE 5 : Presque fini ===
    \node[font=\scriptsize\bfseries, ucablue] at (8.5, -1.5) {5. Pioche};
    \foreach \x in {0,...,4} \foreach \y in {-6,...,-2} \fill[black!80] (6+\x, \y) rectangle ++(1,1);
    \fill[white] (7,-3) rectangle ++(1,1);
    \fill[white] (8,-3) rectangle ++(1,1);
    \fill[white] (9,-3) rectangle ++(1,1);
    \fill[white] (7,-4) rectangle ++(1,1);
    \fill[white] (7,-5) rectangle ++(1,1);
    \fill[white] (9,-4) rectangle ++(1,1);
    \fill[white] (9,-5) rectangle ++(1,1);
    \fill[white] (8,-5) rectangle ++(1,1);
    \fill[orange!50] (7,-7) rectangle ++(1,1);
    
    % === ÉTAPE 6 : Terminé ===
    \node[font=\scriptsize\bfseries, vertok] at (14.5, -1.5) {6. Fini};
    \foreach \x in {0,...,4} \foreach \y in {-6,...,-2} \fill[black!80] (12+\x, \y) rectangle ++(1,1);
    \fill[white] (13,-3) rectangle ++(1,1);
    \fill[white] (14,-3) rectangle ++(1,1);
    \fill[white] (15,-3) rectangle ++(1,1);
    \fill[white] (13,-4) rectangle ++(1,1);
    \fill[white] (13,-5) rectangle ++(1,1);
    \fill[white] (15,-4) rectangle ++(1,1);
    \fill[white] (15,-5) rectangle ++(1,1);
    \fill[white] (14,-5) rectangle ++(1,1);
    \fill[vertok!30] (13,-3) rectangle ++(1,1);
    \fill[vertok!30] (15,-5) rectangle ++(1,1);
    
    % Légende
    \node[right, font=\scriptsize] at (17.5, 3) {\colorbox{yellow!50}{~S~} départ};
    \node[right, font=\scriptsize] at (17.5, 2) {\colorbox{orange!50}{~?~} liste};
    \node[right, font=\scriptsize] at (17.5, 1) {\colorbox{white}{~~} couloir};
    \node[right, font=\scriptsize] at (17.5, 0) {\colorbox{black!80}{\textcolor{white}{~~}} mur};
\end{tikzpicture}

\vspace{0.2cm}
\small \textcolor{ucablue}{\faRandom} Pioche AU HASARD dans la \textbf{liste} ~$\neq$~ DFS qui suit le dernier point

\end{frame}

%% --------------------------------------------------------------------
\begin{frame}{Benchmark : DFS vs Prim}

\centering
\begin{tikzpicture}[scale=0.45]
    % ========== DFS : serpentin pur (1 seul chemin, 0 intersection) ==========
    \node[font=\bfseries, ucablue] at (3.5, 8.5) {DFS};
    \fill[black!80] (0,0) rectangle (7,8);
    % Serpentin : colonne 1 monte, tourne, colonne 3 descend, tourne, colonne 5 monte
    \foreach \x/\y in {1/1, 1/2, 1/3, 1/4, 1/5, 1/6,
                        2/6,
                        3/6, 3/5, 3/4, 3/3, 3/2, 3/1,
                        4/1,
                        5/1, 5/2, 5/3, 5/4, 5/5, 5/6}{
        \fill[white] (\x, \y) rectangle ++(1,1);
    }
    % Culs-de-sac (les 2 bouts du serpentin)
    \fill[red!30] (1,1) rectangle ++(1,1);
    \fill[red!30] (5,6) rectangle ++(1,1);
    % Flèche montrant le chemin unique
    \draw[-{Stealth}, thick, red!40, dashed] (1.5,1.5) -- (1.5,6.5) -- (3.5,6.5) -- (3.5,1.5) -- (5.5,1.5) -- (5.5,6.5);

    % ========== Prim : grille ouverte (5 intersections, 0 cul-de-sac) ==========
    \node[font=\bfseries, ucablue] at (12.5, 8.5) {Prim};
    \fill[black!80] (9,0) rectangle (16,8);
    % Structure en grille : 3 colonnes + 3 rangées connectées
    \foreach \x/\y in {10/1, 11/1, 12/1, 13/1, 14/1,
                        10/2, 12/2, 14/2,
                        10/3, 11/3, 12/3, 13/3, 14/3,
                        10/4, 12/4, 14/4,
                        10/5, 12/5, 14/5,
                        10/6, 11/6, 12/6, 13/6, 14/6}{
        \fill[white] (\x, \y) rectangle ++(1,1);
    }
    % Intersections (3+ voisins) marquées en vert
    \fill[green!30] (10,3) rectangle ++(1,1);  % 3 voisins
    \fill[green!30] (12,1) rectangle ++(1,1);  % 3 voisins
    \fill[green!30] (12,3) rectangle ++(1,1);  % 4 voisins !
    \fill[green!30] (14,3) rectangle ++(1,1);  % 3 voisins
    \fill[green!30] (12,6) rectangle ++(1,1);  % 3 voisins

    % Tableau comparatif
    \node[font=\small] at (3.5, -1.5) {\textcolor{rougeko}{\faTimes} 2 culs-de-sac};
    \node[font=\small] at (3.5, -2.3) {\textcolor{rougeko}{\faTimes} 0 intersection};
    \node[font=\small] at (12.5, -1.5) {\textcolor{vertok}{\faCheck} 0 cul-de-sac};
    \node[font=\small] at (12.5, -2.3) {\textcolor{vertok}{\faCheck} 5 intersections};

    % Flèche verdict
    \draw[-{Stealth}, ultra thick, vertok] (17, 4) -- (18.5, 4);
    \node[font=\bfseries, vertok, right] at (18.5, 4) {Prim};
    \node[font=\small, right] at (18.5, 3) {retenu};
\end{tikzpicture}

\end{frame}

%% --------------------------------------------------------------------
\begin{frame}{Ajout de boucles -- \texttt{add\_loops(maze, loop\_percent)}}

\centering
\begin{tikzpicture}[scale=0.42]
    %% ============================================================
    %% Grille Prim 9x9  --  Cellules impaires (1,3,5,7)
    %% Arbre couvrant : colonnes verticales + 3 ponts horizontaux
    %% Murs candidats = "?" blanc sur fond noir (pas de couleur)
    %% Murs pétés = case blanche + petit check noir
    %% ============================================================

    %% === 0\% BOUCLES (parfait) ===
    \node[font=\small\bfseries, ucablue] at (4.5, 10) {0\% boucles (parfait)};
    \fill[black!80] (0,0) rectangle (9,9);
    % Couloirs : cellules + passages de l'arbre couvrant
    \foreach \x/\y in {
        1/1, 1/2, 1/3, 1/4, 1/5, 1/6, 1/7,
        3/1, 3/2, 3/3, 3/4, 3/5, 3/6, 3/7,
        5/1, 5/2, 5/3, 5/4, 5/5, 5/6, 5/7,
        7/1, 7/2, 7/3, 7/4, 7/5, 7/6, 7/7,
        2/7, 4/5, 6/7}{
        \fill[white] (\x, \y) rectangle ++(1,1);
    }
    % 9 murs candidats : juste un "?" blanc sur le mur noir
    \foreach \px/\py in {2/1,2/3,2/5,4/1,4/3,4/7,6/1,6/3,6/5}{
        \node[font=\tiny, white] at (\px+0.5, \py+0.5) {\textbf{?}};
    }
    \node[font=\scriptsize, gray] at (4.5, -0.8) {9 murs candidats};
    \node[font=\scriptsize, rougeko] at (4.5, -1.6) {1 seul chemin entre 2 cases};

    %% Flèche 25\%
    \draw[-{Stealth}, ultra thick, vertok] (9.8, 4.5) -- (11.3, 4.5);
    \node[font=\scriptsize\bfseries, vertok] at (10.55, 5.5) {25\%};
    \node[font=\tiny, gray] at (10.55, 3.5) {$\lfloor 9 \times 0.25 \rfloor$};
    \node[font=\tiny, gray] at (10.55, 3) {= 2 murs};

    %% === 25\% BOUCLES ===
    \node[font=\small\bfseries, vertok] at (15.5, 10) {25\% boucles \textcolor{ucablue}{(Pac-Man)}};
    \fill[black!80] (11,0) rectangle (20,9);
    \foreach \x/\y in {
        12/1, 12/2, 12/3, 12/4, 12/5, 12/6, 12/7,
        14/1, 14/2, 14/3, 14/4, 14/5, 14/6, 14/7,
        16/1, 16/2, 16/3, 16/4, 16/5, 16/6, 16/7,
        18/1, 18/2, 18/3, 18/4, 18/5, 18/6, 18/7,
        13/7, 15/5, 17/7}{
        \fill[white] (\x, \y) rectangle ++(1,1);
    }
    % 2 murs pétés → deviennent blancs + check
    \fill[white] (15,3) rectangle ++(1,1);
    \node[font=\tiny] at (15.5, 3.5) {\checkmark};
    \fill[white] (17,1) rectangle ++(1,1);
    \node[font=\tiny] at (17.5, 1.5) {\checkmark};
    % Murs candidats restants : "?" blanc sur noir
    \foreach \px/\py in {13/1,13/3,13/5,15/1,15/7,17/3,17/5}{
        \node[font=\tiny, white] at (\px+0.5, \py+0.5) {\textbf{?}};
    }
    \node[font=\scriptsize, vertok] at (15.5, -0.8) {2 boucles créées};
    \node[font=\scriptsize, vertok] at (15.5, -1.6) {choix de chemins \faArrowRight~jouable};

    %% Flèche 50\%
    \draw[-{Stealth}, ultra thick, orangewarn] (20.8, 4.5) -- (22.3, 4.5);
    \node[font=\scriptsize\bfseries, orangewarn] at (21.55, 5.5) {50\%};
    \node[font=\tiny, gray] at (21.55, 3.5) {$\lfloor 9 \times 0.5 \rfloor$};
    \node[font=\tiny, gray] at (21.55, 3) {= 4 murs};

    %% === 50\% BOUCLES (trop ouvert) ===
    \node[font=\small\bfseries, orangewarn] at (26.5, 10) {50\% boucles (trop ouvert)};
    \fill[black!80] (22,0) rectangle (31,9);
    \foreach \x/\y in {
        23/1, 23/2, 23/3, 23/4, 23/5, 23/6, 23/7,
        25/1, 25/2, 25/3, 25/4, 25/5, 25/6, 25/7,
        27/1, 27/2, 27/3, 27/4, 27/5, 27/6, 27/7,
        29/1, 29/2, 29/3, 29/4, 29/5, 29/6, 29/7,
        24/7, 26/5, 28/7}{
        \fill[white] (\x, \y) rectangle ++(1,1);
    }
    % 4 murs pétés → blancs + check
    \fill[white] (26,3) rectangle ++(1,1);
    \node[font=\tiny] at (26.5, 3.5) {\checkmark};
    \fill[white] (28,1) rectangle ++(1,1);
    \node[font=\tiny] at (28.5, 1.5) {\checkmark};
    \fill[white] (24,5) rectangle ++(1,1);
    \node[font=\tiny] at (24.5, 5.5) {\checkmark};
    \fill[white] (28,5) rectangle ++(1,1);
    \node[font=\tiny] at (28.5, 5.5) {\checkmark};
    % Murs candidats restants : "?" blanc sur noir
    \foreach \px/\py in {24/1,24/3,26/1,26/7,28/3}{
        \node[font=\tiny, white] at (\px+0.5, \py+0.5) {\textbf{?}};
    }
    \node[font=\scriptsize, orangewarn] at (26.5, -0.8) {4 boucles --- trop de raccourcis};
    \node[font=\scriptsize, rougeko] at (26.5, -1.6) {le labyrinthe perd son intérêt};

    %% Légende
    \node[font=\scriptsize] at (4.5, -3) {\textcolor{white}{\textbf{?}}~sur mur = mur candidat};
    \node[font=\scriptsize] at (15.5, -3) {\checkmark~= mur supprimé (boucle)};
    \node[font=\scriptsize] at (26.5, -3) {noir = mur ~\textbar~ blanc = couloir};
\end{tikzpicture}

\vspace{0.1cm}
\small \textbf{Pac-Man} $\approx$ \textcolor{vertok}{25\%} ~\textbar~ Trop facile $>$ 50\% ~\textbar~ \texttt{add\_loops(maze, \%)} supprime un \% des murs candidats

\end{frame}

%% ####################################################################
%% PARTIE NADA
%% ####################################################################

\section{Interface Graphique Pygame}

%% --------------------------------------------------------------------
\begin{frame}{Architecture globale}

\centering
\begin{tikzpicture}[
    box/.style={draw, rounded corners, minimum width=2.5cm, minimum height=1.5cm, align=center, font=\small, thick},
    arr/.style={-{Stealth[length=3mm]}, very thick, ucablue}
]
    % Client
    \node[box, fill=gray!10] (client) at (0, 0) {\faDesktop\\[0.1cm]\textbf{Client}\\navigateur / code};
    
    % Internet
    \node[font=\Large] (net) at (3.5, 0) {\faGlobe};
    \node[font=\scriptsize, below=0.1cm of net] {HTTP};
    
    % Render
    \node[box, fill=green!10] (render) at (7, 0) {\faCloud\\[0.1cm]\textbf{Render}\\gunicorn};
    
    % Flask
    \node[box, fill=ucalight] (flask) at (11, 0) {\faPython\\[0.1cm]\textbf{Flask}\\app.py};
    
    % Générateur
    \node[box, fill=orange!10] (gen) at (11, -2.5) {\faDice\\[0.1cm]\textbf{Prim}\\maze\_Prim\_loops};
    
    % JSON retour
    \node[box, fill=yellow!10] (json) at (3.5, -2.5) {\faFileCode\\[0.1cm]\textbf{JSON}\\maze + méta};
    
    % Flèches
    \draw[arr] (client) -- node[above, font=\scriptsize] {GET /maze} (net);
    \draw[arr] (net) -- (render);
    \draw[arr] (sqlite) -- (pygame);
    \draw[arr] (pygame) -- (gen);
    \draw[arr, orange] (gen) -| (json);
    \draw[arr, orange] (json) -- node[below, font=\scriptsize] {réponse} (client);
    
\end{tikzpicture}

\vspace{0.3cm}
\small URL : \textcolor{ucablue}{\texttt{https://ter-s2-d.onrender.com/maze?width=21\&height=21\&loop\_percent=25}}

\end{frame}

%% --------------------------------------------------------------------
\begin{frame}[fragile]{API : Routes \& Réponse JSON}

\begin{columns}[T]
\begin{column}{0.48\textwidth}
\begin{block}{\small Routes}
\small
\begin{tabular}{ll}
\texttt{GET /} & Health check \\
\texttt{GET /maze} & Générer un maze \\
\end{tabular}

\vspace{0.3cm}
\textbf{Paramètres de \texttt{/maze}} :
\begin{tabular}{lll}
\texttt{width} & int & défaut: 21 \\
\texttt{height} & int & défaut: 21 \\
\texttt{loop\_percent} & int & défaut: 25 \\
\end{tabular}
\end{block}
\end{column}

\begin{column}{0.48\textwidth}
\begin{block}{\small Réponse JSON}
\begin{lstlisting}[language=Python]
{
  "width": 21,
  "height": 21,
  "loop_percent": 25,
  "maze": [
    [0, 0, 0, 0, ...],
    [0, 1, 1, 0, ...],
    ...
  ]
}
\end{lstlisting}
\end{block}
\small \textbf{Convention} : 0 = mur, 1 = couloir
\end{column}
\end{columns}

\end{frame}

%% --------------------------------------------------------------------
\begin{frame}{Déploiement Render \& Deploy Hook}

\begin{columns}[T]
\begin{column}{0.55\textwidth}
\centering
\begin{tikzpicture}[
    box/.style={draw, rounded corners, minimum width=2.5cm, minimum height=1cm, align=center, font=\scriptsize, thick},
    arr/.style={-{Stealth[length=2mm]}, thick}
]
    % GitHub IkramB23
    \node[box, fill=gray!10] (ikram) at (0, 3) {\faGithub~\textbf{IkramB23}\\TER\_S2\_D};
    
    % CI
    \node[box, fill=ucalight] (ci) at (0, 1.5) {\faCog~\textbf{GitHub Actions}\\test + deploy};
    
    % Deploy Hook
    \node[box, fill=yellow!15] (hook) at (4, 1.5) {\faLink~\textbf{Deploy Hook}\\curl (secret)};
    
    % GitHub webNada
    \node[box, fill=gray!10] (nada) at (4, 3) {\faGithub~\textbf{webNada}\\TER\_S2\_D};
    
    % Render
    \node[box, fill=green!15] (render) at (4, 0) {\faCloud~\textbf{Render}\\gunicorn app:app};
    
    % Flèches
    \draw[arr, ucablue] (ikram) -- node[left, font=\scriptsize] {push} (ci);
    \draw[arr, vertok] (ci) -- node[above, font=\scriptsize] {si \textcolor{vertok}{\faCheck}} (hook);
    \draw[arr, orange] (hook) -- (render);
    \draw[arr, gray, dashed] (ikram) -- node[above, font=\scriptsize] {sync} (nada);
    \draw[arr, gray] (nada) -- (render);
    
\end{tikzpicture}
\end{column}

\begin{column}{0.42\textwidth}
\begin{block}{\small Stack technique}
\small
\begin{tabular}{ll}
\faPython & Python / Flask \\
\faServer & Gunicorn (prod) \\
\faCloud & Render (gratuit) \\
\faGithub & 2 repos synchro \\
\faLock & Secret GitHub \\
\end{tabular}
\end{block}

\vspace{0.2cm}
\begin{alertblock}{\scriptsize Note}
\scriptsize Le serveur Render ``dort'' après 15 min\\
Premier appel $\sim$30s de réveil
\end{alertblock}
\end{column}
\end{columns}

\end{frame}

%% ====================================================================
\section{Tests}

%% --------------------------------------------------------------------
\begin{frame}{Stratégie de tests : 4 niveaux}

\centering
\begin{tikzpicture}[
    testlevel/.style={draw, rounded corners, minimum width=3.5cm, minimum height=1cm, font=\small, align=center},
    arr/.style={-{Stealth[length=2mm]}, thick, gray}
]
    % 4 niveaux
    \node[testlevel, fill=ucalight] (func) at (0, 0) {\textbf{Tests fonctionnels}\\{\scriptsize dimensions, bords, valeurs}};
    \node[testlevel, fill=green!10] (struct) at (5.5, 0) {\textbf{Tests structurels}\\{\scriptsize connexité BFS, culs-de-sac}};
    \node[testlevel, fill=orange!10] (api) at (0, -2.5) {\textbf{Tests API locale}\\{\scriptsize routes, JSON, paramètres}};
    \node[testlevel, fill=red!10] (e2e) at (5.5, -2.5) {\textbf{Tests bout en bout}\\{\scriptsize API sur Render (cloud)}};
    
    % Flèches
    \draw[arr] (func) -- (struct);
    \draw[arr] (func) -- (api);
    \draw[arr] (api) -- (e2e);
    \draw[arr] (struct) -- (e2e);
    
    % Compteurs
    \node[font=\scriptsize, ucablue, above=0.1cm of func] {\faCheckCircle~6 tests (Ikram)};
    \node[font=\scriptsize, ucablue, above=0.1cm of struct] {\faCheckCircle~4 tests (Ikram)};
    \node[font=\scriptsize, ucablue, below=0.1cm of api] {\faCheckCircle~5 tests (Ikram)};
    \node[font=\scriptsize, ucablue, below=0.1cm of e2e] {\faCheckCircle~tests (Nada)};
    
    % Total
    \node[draw=vertok, thick, rounded corners, fill=green!10, font=\bfseries] at (2.75, -4.5) {\textcolor{vertok}{\faCheck}~15 tests dans le CI ~\textbar~ pytest};
\end{tikzpicture}

\end{frame}

%% --------------------------------------------------------------------
\begin{frame}[fragile]{Exemples de tests clés}

\begin{columns}[T]
\begin{column}{0.48\textwidth}
\begin{block}{\scriptsize Test fonctionnel : dimensions}
\begin{lstlisting}[language=Python]
def test_dimensions():
    maze = generate_pacman_maze(21, 21)
    assert len(maze) == 21
    assert len(maze[0]) == 21
\end{lstlisting}
\end{block}

\begin{block}{\scriptsize Test structural : connexité BFS}
\begin{lstlisting}[language=Python]
def test_connectivity():
    # BFS depuis un couloir
    # verifie qu'on atteint TOUS
    # les couloirs du labyrinthe
    visited = bfs(maze, start)
    assert len(visited) == total
\end{lstlisting}
\end{block}
\end{column}

\begin{column}{0.48\textwidth}
\begin{block}{\scriptsize Test API locale}
\begin{lstlisting}[language=Python]
def test_maze_route():
    client = app.test_client()
    resp = client.get('/maze')
    assert resp.status_code == 200
    data = resp.get_json()
    assert data["width"] == 21
\end{lstlisting}
\end{block}

\begin{block}{\scriptsize Test bout en bout (Nada)}
\begin{lstlisting}[language=Python]
def test_cloud():
    URL = "https://ter-s2-d.
           onrender.com/maze"
    resp = requests.get(URL)
    assert resp.status_code == 200
\end{lstlisting}
\end{block}
\end{column}
\end{columns}

\end{frame}

%% ====================================================================
\section{CI/CD}

%% --------------------------------------------------------------------
\begin{frame}{Pipeline CI/CD -- GitHub Actions}

\centering
\begin{tikzpicture}[scale=0.9, every node/.style={transform shape},
    box/.style={draw, rounded corners, minimum width=2cm, minimum height=1.1cm, align=center, font=\small, thick},
    arr/.style={-{Stealth[length=3mm]}, very thick},
    note/.style={font=\scriptsize, text=gray, align=center}
]
    % Trigger
    \node[box, fill=gray!20] (push) at (0, 0) {\faCodeBranch\\[0.1cm]\texttt{git push}\\sur \texttt{main}};
    
    % Job Test
    \node[box, fill=ucalight] (test) at (3.5, 0) {\faCog\\[0.1cm]\textbf{Job: test}\\15 tests};
    
    % Décision
    \node[diamond, draw, thick, fill=yellow!20, minimum size=1.1cm, inner sep=0pt, font=\scriptsize] (check) at (6.5, 0) {\textbf{OK?}};
    
    % Job Deploy
    \node[box, fill=green!15] (deploy) at (9.5, 0) {\faCloud\\[0.1cm]\textbf{deploy}\\curl hook};
    
    % Render
    \node[box, fill=vertok!20] (render) at (12.5, 0) {\faServer\\[0.1cm]\textbf{Render}\\redéploie};
    
    % Flèches
    \draw[arr, ucablue] (push) -- (test);
    \draw[arr, ucablue] (test) -- (check);
    \draw[arr, vertok] (check) -- node[above, font=\scriptsize, vertok] {\faCheck} (deploy);
    \draw[arr, vertok] (deploy) -- (render);
    
    % Flèche rouge (échec)
    \draw[arr, rougeko] (check) -- ++(0, -1.6) node[below, font=\scriptsize, rougeko] {\faTimes~\textbf{STOP}};
    
    % Annotations
    \node[note, above=0.4cm of test] {Ubuntu\\Python 3.10};
    \node[note, above=0.4cm of deploy] {\texttt{needs: test}};
    \node[note, above=0.4cm of render] {ter-s2-d};

\end{tikzpicture}

\vspace{0.4cm}
\small \textcolor{ucablue}{\faFileCode~\texttt{.github/workflows/ci.yml}} -- déclenché automatiquement à chaque push

\end{frame}

%% --------------------------------------------------------------------
\begin{frame}{Démo rouge / vert}

\begin{columns}[c]
\begin{column}{0.48\textwidth}
\centering
\begin{tikzpicture}[
    stepbox/.style={draw, rounded corners, minimum width=4cm, minimum height=1.5cm, align=center, font=\small, thick}
]
    % Étape 1 : casser
    \node[stepbox, fill=red!10, draw=rougeko] (break) at (0, 2) {
        \textcolor{rougeko}{\faTimesCircle~\textbf{On casse exprès}}\\[0.1cm]
        \scriptsize\texttt{assert len(maze) == 99}\\
        \scriptsize (au lieu de 21)
    };
    
    % Push
    \node[font=\scriptsize] (push1) at (0, 0.5) {\texttt{git push} $\rightarrow$ CI lance les tests};
    
    % Résultat rouge
    \node[stepbox, fill=red!20, draw=rougeko] (red) at (0, -1) {
        \textcolor{rougeko}{\faTimes~\textbf{CI ROUGE}}\\
        \scriptsize 1 failed, 14 passed\\
        \scriptsize \textbf{Pas de déploiement}
    };
    
    \draw[-{Stealth}, thick, rougeko] (break) -- (push1);
    \draw[-{Stealth}, thick, rougeko] (push1) -- (red);
\end{tikzpicture}
\end{column}

\begin{column}{0.48\textwidth}
\centering
\begin{tikzpicture}[
    stepbox/.style={draw, rounded corners, minimum width=4cm, minimum height=1.5cm, align=center, font=\small, thick}
]
    % Étape 2 : réparer
    \node[stepbox, fill=green!10, draw=vertok] (fix) at (0, 2) {
        \textcolor{vertok}{\faCheckCircle~\textbf{On répare}}\\[0.1cm]
        \scriptsize\texttt{assert len(maze) == 21}\\
        \scriptsize (valeur correcte)
    };
    
    % Push
    \node[font=\scriptsize] (push2) at (0, 0.5) {\texttt{git push} $\rightarrow$ CI relance};
    
    % Résultat vert
    \node[stepbox, fill=green!20, draw=vertok] (green) at (0, -1) {
        \textcolor{vertok}{\faCheck~\textbf{CI VERT}}\\
        \scriptsize 15 passed\\
        \scriptsize \textbf{Déploiement déclenché}
    };
    
    \draw[-{Stealth}, thick, vertok] (fix) -- (push2);
    \draw[-{Stealth}, thick, vertok] (push2) -- (green);
\end{tikzpicture}
\end{column}
\end{columns}

\vspace{0.4cm}
\centering
\small \textcolor{ucablue}{\faShieldAlt}~Le CI \textbf{protège main} : code cassé $\neq$ déployé en production

\end{frame}

%% ####################################################################
%% PARTIE AYA -- PLACEHOLDER
%% ####################################################################

\section{Interface graphique}

%% --------------------------------------------------------------------
\begin{frame}{Interface graphique -- Aya}
\centering
\vspace{2cm}
{\Large\textcolor{gray}{Interface graphique}}\\[0.5cm]
{\large\textcolor{gray}{Aya Haddoun}}\\[1cm]
\textcolor{gray!50}{\scriptsize (slides à compléter par Aya)}
\end{frame}

%% --------------------------------------------------------------------
\begin{frame}{Interface graphique -- Aya}
\centering
\vspace{2cm}
{\Large\textcolor{gray}{Interface graphique}}\\[0.5cm]
{\large\textcolor{gray}{Aya Haddoun}}\\[1cm]
\textcolor{gray!50}{\scriptsize (slides à compléter par Aya)}
\end{frame}

%% --------------------------------------------------------------------
\begin{frame}{Interface graphique -- Aya}
\centering
\vspace{2cm}
{\Large\textcolor{gray}{Interface graphique}}\\[0.5cm]
{\large\textcolor{gray}{Aya Haddoun}}\\[1cm]
\textcolor{gray!50}{\scriptsize (slides à compléter par Aya)}
\end{frame}

%% ####################################################################
%% PARTIE COMMUNE : Structure & Collaboration
%% ####################################################################

\section{Projet \& Collaboration}

%% --------------------------------------------------------------------
\begin{frame}{Structure du projet}

\begin{columns}[T]
\begin{column}{0.45\textwidth}
\centering
\begin{tikzpicture}[
    folder/.style={draw, thick, rounded corners, fill=ucalight, minimum width=2cm, font=\small},
    file/.style={font=\scriptsize\ttfamily, anchor=west},
    edge from parent/.style={draw, thick},
    level 1/.style={sibling distance=3.5cm, level distance=1.2cm},
    grow=down
]
    \node[folder] {\faFolder~TER\_S2\_D}
    [edge from parent]
    child { node[folder] {\faFolder~Proto}
        child { node[file, yshift=-0.2cm] {
            \begin{minipage}{3cm}
            \faFile~maze\_dfs.py\\
            \faFile~maze\_Prim.py\\
            \faFile~maze\_Prim\_loops.py\\
            \faFile~maze\_analyzer.py\\
            \faFile~Benchmark\_prim.py
            \end{minipage}
        }}
    }
    child { node[folder] {\faFolder~tests}
        child { node[file, yshift=-0.2cm] {
            \begin{minipage}{3.2cm}
            \faFile~test\_maze\_generator.py\\
            \faFile~test\_app\_local.py
            \end{minipage}
        }}
    };
\end{tikzpicture}
\end{column}

\begin{column}{0.50\textwidth}
\centering
\begin{tikzpicture}[
    file/.style={font=\scriptsize\ttfamily, anchor=west}
]
    \node[file] at (0, 4.5) {\faFile~local_pacman.py \textcolor{gray}{\tiny (API Flask)}};
    \node[file] at (0, 4.0) {\faFile~requirements.txt};
    \node[file] at (0, 3.5) {\faFile~.gitignore};
    \node[file] at (0, 3.0) {\faFile~test\_api\_cloud.py \textcolor{gray}{\tiny (Nada)}};
    \node[file] at (0, 2.5) {\faFile~test\_Lab\_properties.py \textcolor{gray}{\tiny (Nada)}};
    \node[file] at (0, 2.0) {\faFileImage~Gallery/best\_ascii.png};
    \node[file] at (0, 1.5) {\faFile~.github/workflows/ci.yml};
\end{tikzpicture}

\vspace{0.3cm}
\begin{block}{\small Choix techniques}
\scriptsize
\begin{tabular}{ll}
Prim (pas DFS) & plus adapté Pac-Man \\
Fonctions (pas classes) & plus testable \\
Pygame & bibliothèque 2D légère \\
Render (gratuit) & déploiement facile \\
Dimensions impaires & structure de grille \\
\end{tabular}
\end{block}
\end{column}
\end{columns}

\end{frame}

%% --------------------------------------------------------------------
\begin{frame}{Répartition \& Collaboration}

\centering
\begin{tikzpicture}[
    member/.style={draw, rounded corners, minimum width=3cm, minimum height=1cm, align=center, font=\small, thick},
    task/.style={font=\scriptsize, anchor=west},
    arr/.style={-{Stealth[length=2mm]}, thick, gray}
]
    % Ikram
    \node[member, fill=blue!10] (ikram) at (0, 2) {\faUser\\[0.05cm]\textbf{Ikram}};
    \node[task] at (2, 2.8) {$\bullet$ Algo DFS + adaptation Prim};
    \node[task] at (2, 2.3) {$\bullet$ Benchmark comparatif};
    \node[task] at (2, 1.8) {$\bullet$ 15 tests (pytest)};
    \node[task] at (2, 1.3) {$\bullet$ CI/CD (GitHub Actions)};
    \node[task] at (2, 0.8) {$\bullet$ Gallery + .gitignore + sync repos};

    
    % Nada
    \node[member, fill=green!10] (nada) at (0, -0.5) {\faUser\\[0.05cm]\textbf{Nada}};
    \node[task] at (2, 0) {$\bullet$ Algo Prim (original, classe)};
    \node[task] at (2, -0.5) {$\bullet$ Aperçu Pygame + app.py};
    \node[task] at (2, -1) {$\bullet$ Déploiement Render};
    \node[task] at (2, -1.5) {$\bullet$ Tests cloud (end-to-end)};
    
    % Aya
    \node[member, fill=orange!10] (aya) at (0, -3) {\faUser\\[0.05cm]\textbf{Aya}};
    \node[task] at (2, -2.7) {$\bullet$ Interface graphique (GUI)};
    \node[task] at (2, -3.2) {$\bullet$ Appel API + affichage visuel};
    
    % Git
    \node[draw, dashed, rounded corners, fill=gray!5, minimum width=4cm, font=\scriptsize, align=center] at (10, 0) {
        \faGithub~\textbf{GitHub}\\[0.2cm]
        Repo principal : IkramB23\\
        Fork Nada : webNada\\
        Branches : main, ikram-work, nada\_work\\
        Merge sur main\\
        Prof : gmenez (collaborateur)
    };
\end{tikzpicture}

\end{frame}

%% ====================================================================
%% RÉFÉRENCES
%% ====================================================================

\section*{Références}

\begin{frame}{Références}
\scriptsize
\setlength{\itemsep}{0pt}
\begin{thebibliography}{9}
    \bibitem{render} Render -- Your First Deploy ~~\url{https://render.com/docs/your-first-deploy}
    \bibitem{flask} Pygame -- Documentation officielle ~~\\url{https://www.pygame.org/docs/}
    \bibitem{realpython} Real Python -- API Integration ~~\url{https://realpython.com/api-integration-in-python/}
    \bibitem{pytest} pytest -- Documentation ~~\url{https://docs.pytest.org/}
    \bibitem{ghactions} GitHub Actions -- Documentation ~~\url{https://docs.github.com/en/actions}
    \bibitem{prim} Algorithme de Prim -- Wikipédia ~~\url{https://fr.wikipedia.org/wiki/Algorithme_de_Prim}
    \bibitem{dfs} DFS Maze Generation ~~\url{https://en.wikipedia.org/wiki/Maze_generation_algorithm}
\end{thebibliography}
\end{frame}

%% ====================================================================
\begin{frame}{}
\centering
\vspace{2cm}
{\Huge\textcolor{ucablue}{Merci !}}\\[1cm]
{\large Questions ?}\\[0.5cm]
{\small\textcolor{gray}{\faGithub~github.com/IkramB23/TER\_S2\_D}}\\
{\small\textcolor{gray}{\faGlobe~ter-s2-d.onrender.com}}
\end{frame}

\end{document}
